\documentclass[a4paper,14pt]{extarticle}

% ============ Пакеты ============
\usepackage{fontspec}
\usepackage[russian,english]{babel}

\setmainfont{Times New Roman}
\setsansfont{Arial}
\setmonofont{Courier New}
\usepackage{geometry}
\geometry{left=2.5cm,right=2cm,top=2cm,bottom=2cm}
\usepackage{amsmath,amssymb}
\usepackage{graphicx}
\usepackage{booktabs}
\usepackage{array}
\usepackage{longtable}
\usepackage{listings}
\usepackage{xcolor}
\usepackage{hyperref}
\usepackage{indentfirst}
\usepackage{setspace}
\usepackage{titlesec}
\usepackage{caption}
\usepackage{float}
\usepackage{enumitem}
\usepackage{tikz}
\usetikzlibrary{arrows.meta,positioning}

\onehalfspacing

% ============ Оформление листингов ============
\definecolor{codebg}{RGB}{248,248,248}
\definecolor{codegreen}{RGB}{0,128,0}
\definecolor{codeblue}{RGB}{0,0,180}
\definecolor{codered}{RGB}{180,0,0}
\lstdefinestyle{jsstyle}{
  backgroundcolor=\color{codebg},
  basicstyle=\ttfamily\small,
  keywordstyle=\color{codeblue}\bfseries,
  commentstyle=\color{codegreen}\itshape,
  stringstyle=\color{codered},
  numbers=left,
  numberstyle=\tiny\color{gray},
  numbersep=6pt,
  frame=single,
  rulecolor=\color{gray!40},
  breaklines=true,
  tabsize=2,
  showstringspaces=false,
  language=Java,
  captionpos=b
}
\lstset{style=jsstyle}

\hypersetup{
  colorlinks=true,
  linkcolor=black,
  urlcolor=blue,
  citecolor=black
}

% ============ Титульный лист ============
\title{
  \vspace{-2cm}
  \textbf{Лабораторная работа №3}\\[0.3cm]
  \textbf{Базовые растровые алгоритмы}\\[0.6cm]
  \large Отчёт
}
\author{Студент: Герасимович Данила\\
        Группа: 4, ФМиИС}
\date{Октябрь 2026}

% ======================================
\begin{document}
\maketitle

\tableofcontents
\newpage

% ======================================
\section{Введение}

\subsection{Цель работы}
Закрепление теоретического материала и практическое освоение базовых алгоритмов растеризации отрезков и кривых:
\begin{itemize}
  \item пошаговый алгоритм;
  \item алгоритм ЦДА (цифровой дифференциальный анализатор);
  \item алгоритм Брезенхема для отрезка;
  \item алгоритм Брезенхема для окружности.
\end{itemize}

\subsection{Постановка задачи}
Требуется разработать веб-приложение, иллюстрирующее работу перечисленных алгоритмов, с удобным интерфейсом, сеткой, осями координат, масштабированием, анимацией построения и замером времени выполнения. Также необходимо привести один пример с сопроводительными вычислениями для каждого алгоритма, составить краткий отчёт о временных характеристиках и пояснить, каким образом целочисленные координаты привязаны к дискретной сетке.

\subsection{Среда разработки}
Приложение реализовано в виде единого файла \texttt{index.html}, содержащего HTML-разметку, CSS-стили и JavaScript-код. Запускается в любом современном браузере (Chrome, Firefox, Edge, Opera) без установки дополнительного программного обеспечения.

% ======================================
\section{Теоретическая часть}

\subsection{Понятие растеризации}
Растеризация (сканирующее преобразование) — процесс преобразования геометрического примитива, заданного в непрерывных координатах, в набор пикселей дискретного растра. Качество растеризации определяется близостью получающегося набора пикселей к идеальной математической линии и отсутствием визуально заметных разрывов.

\subsection{Пошаговый алгоритм}
Идея пошагового (прямого) алгоритма состоит в том, что вдоль ведущей оси (той, по которой приращение координаты больше) выполняется шаг ровно на 1 пиксель, а вторая координата вычисляется по уравнению прямой и округляется:
\[
y = y_1 + \frac{dy}{dx}\,(x - x_1), \quad dx = x_2 - x_1,\; dy = y_2 - y_1.
\]
Если $|dx| \ge |dy|$, ведущей является ось $x$; в противном случае — ось $y$. Главный недостаток метода — операция деления с плавающей точкой на каждом шаге.

\subsection{Алгоритм ЦДА (DDA)}
Цифровой дифференциальный анализатор основан на вычислении приращений:
\[
\Delta x = \frac{x_2-x_1}{\mathrm{steps}},\quad \Delta y = \frac{y_2-y_1}{\mathrm{steps}},\quad \mathrm{steps} = \max(|x_2-x_1|,\,|y_2-y_1|).
\]
На каждом шаге координаты увеличиваются на эти приращения и округляются до целых:
\[
x_{i+1} = x_i + \Delta x,\quad y_{i+1} = y_i + \Delta y.
\]
Алгоритм также использует арифметику с плавающей точкой, но требует лишь $O(\max(|dx|,|dy|))$ шагов.

\subsection{Алгоритм Брезенхема для отрезка}
Алгоритм Брезенхема работает только с целочисленной арифметикой. Основная идея — на каждом шаге выбирать из двух соседних пикселей тот, который ближе к идеальной линии. Для отрезка с $|dx| \ge |dy|$ вводятся:
\[
dx = |x_2 - x_1|,\quad dy = -|y_2 - y_1|,\quad err = dx + dy.
\]
На каждом шаге:
\[
e_2 = 2\,err;\quad
\text{если } e_2 \ge dy,\; \text{то } err \mathrel{+}= dy,\; x \mathrel{+}= s_x;\quad
\text{если } e_2 \le dx,\; \text{то } err \mathrel{+}= dx,\; y \mathrel{+}= s_y.
\]
Здесь $s_x = \operatorname{sign}(x_2-x_1)$, $s_y = \operatorname{sign}(y_2-y_1)$. Алгоритм работает за $O(|dx|+|dy|)$ шагов и не использует операций с плавающей точкой.

\subsection{Алгоритм Брезенхема для окружности}
Для окружности радиуса $R$ с центром $(x_c, y_c)$ используется свойство симметрии: достаточно построить один октант (от $x = 0$, $y = R$ до $x = y = R/\sqrt{2}$), после чего остальные точки получаются отражениями:
\[
(\pm x, \pm y),\quad (\pm y, \pm x).
\]
Начальное значение решающей функции: $d = 3 - 2R$. На каждом шаге при $x \to x+1$:
\[
\text{если } d < 0,\; d \mathrel{+}= 4x + 6,\quad
\text{иначе } d \mathrel{+}= 4(x - y) + 10,\; y \mathrel{-}= 1.
\]
Так же, как и для отрезка, все вычисления целочисленные.

% ======================================
\section{Реализация}

\subsection{Архитектура приложения}
Приложение реализовано как одностраничное веб-приложение на чистом HTML/CSS/JavaScript. Холст \texttt{<canvas>} размером $900\times600$ пикселей используется для вывода растровых точек и системы координат. Все алгоритмы реализованы в виде отдельных функций, возвращающих массив точек \texttt{\{x, y\}} в мировых координатах; отрисовка выполняется отдельным универсальным методом.

Структура модулей:
\begin{itemize}
  \item \texttt{readParams()} — считывание параметров из формы;
  \item \texttt{worldToScreen(x, y)} — преобразование мировых координат в экранные;
  \item \texttt{drawGrid()} — отрисовка сетки, осей и подписей;
  \item \texttt{stepLine}, \texttt{ddaLine}, \texttt{bresenhamLine}, \texttt{bresenhamCircle} — реализации четырёх алгоритмов;
  \item \texttt{draw()}, \texttt{animate()}, \texttt{timing()} — управляющие функции.
\end{itemize}

\subsection{Привязка координат к дискретной сетке}
Пусть масштаб равен $s$ пикселей на единицу мировой координаты, холст имеет размер $W\times H$, начало координат находится в центре холста. Тогда мировая точка $(x, y)$ переводится в экранную:
\[
s_x = \frac{W}{2} + x\,s,\quad s_y = \frac{H}{2} - y\,s.
\]
Знак минус перед $y$ соответствует направлению оси $Y$ вверх (в декартовой системе), тогда как на экране ось $Y$ направлена вниз. Каждая целочисленная точка $(x, y)$ изображается квадратом со стороной $s$ и центром в $(s_x, s_y)$, что визуально соответствует одной ячейке сетки. Значения $x$ и $y$ пробегают целочисленные узлы, поэтому координаты дискретны и однозначно привязаны к сетке.

\subsection{Пользовательский интерфейс}
Интерфейс включает:
\begin{itemize}
  \item выпадающий список выбора алгоритма;
  \item поля ввода $x_1, y_1, x_2, y_2$ (или центра и радиуса для окружности);
  \item регулятор масштаба (10–80 пикселей на единицу);
  \item флажок отображения сетки;
  \item кнопки: \textit{Построить}, \textit{Анимация}, \textit{Замер времени}, \textit{Очистить};
  \item информационную панель с количеством точек;
  \item таблицу растровых точек с подсветкой текущего шага при анимации.
\end{itemize}

% ======================================
\section{Тестирование и примеры вычислений}

Для демонстрации корректности все четыре алгоритма проверены на одинаковых входных данных: отрезок $(2,3) \to (11,7)$ и окружность с центром $(0,0)$ и радиусом $5$.

\subsection{Пошаговый алгоритм}
Дано: $x_1=2,\,y_1=3,\,x_2=11,\,y_2=7$, $dx=9$, $dy=4$. Так как $|dx|>|dy|$, шагаем по $x$. Значение $y$ вычисляется по формуле $y = 3 + \tfrac{4}{9}(x-2)$ и округляется:

\begin{center}
\begin{tabular}{ccccccccccc}
\toprule
$x$ & 2 & 3 & 4 & 5 & 6 & 7 & 8 & 9 & 10 & 11\\
$y_{\text{неокругл.}}$ & 3.00 & 3.44 & 3.89 & 4.33 & 4.78 & 5.22 & 5.67 & 6.11 & 6.56 & 7.00\\
$y_{\text{округл.}}$ & 3 & 3 & 4 & 4 & 5 & 5 & 6 & 6 & 7 & 7\\
\bottomrule
\end{tabular}
\end{center}

Получаем точки:
\[
(2,3), (3,3), (4,4), (5,4), (6,5), (7,5), (8,6), (9,6), (10,7), (11,7).
\]

\subsection{Алгоритм ЦДА}
$\mathrm{steps} = \max(9,4) = 9$, $\Delta x = 1$, $\Delta y = 4/9 \approx 0{,}4444$.
Накопление координат:
\begin{center}
\begin{tabular}{ccccccccccc}
\toprule
$i$ & 0 & 1 & 2 & 3 & 4 & 5 & 6 & 7 & 8 & 9\\
$x$ & 2.00 & 3.00 & 4.00 & 5.00 & 6.00 & 7.00 & 8.00 & 9.00 & 10.00 & 11.00\\
$y$ & 3.000 & 3.444 & 3.889 & 4.333 & 4.778 & 5.222 & 5.667 & 6.111 & 6.556 & 7.000\\
$\mathrm{round}(y)$ & 3 & 3 & 4 & 4 & 5 & 5 & 6 & 6 & 7 & 7\\
\bottomrule
\end{tabular}
\end{center}

Результат совпадает с пошаговым алгоритмом:
\[
(2,3), (3,3), (4,4), (5,4), (6,5), (7,5), (8,6), (9,6), (10,7), (11,7).
\]

\subsection{Алгоритм Брезенхема для отрезка}
$dx = 9$, $dy = -4$, $s_x = 1$, $s_y = 1$, $err = 5$.
Пошаговая таблица:
\begin{center}
\begin{tabular}{cccccl}
\toprule
№ & $x$ & $y$ & $err$ & $e_2$ & Действие \\
\midrule
1 & 2 & 3 & 5 & 10 & $e_2 \ge -4 \Rightarrow x{+}{=}1$; $e_2 \le 9 \Rightarrow y{+}{=}1$ \\
2 & 3 & 4 & 10 & 20 & $x{+}{=}1$, $y{+}{=}1$ \\
$\dots$ & & & & & \\
\bottomrule
\end{tabular}
\end{center}

Получаемые точки:
\[
(2,3), (3,4), (4,4), (5,5), (6,5), (7,6), (8,6), (9,7), (10,7), (11,8).
\]

Заметим, что Брезенхем даёт ту же «лестницу», что и пошаговый алгоритм с ЦДА (разница в выборе соседней точки при равенстве расстояний), но без операций с плавающей точкой.

\subsection{Алгоритм Брезенхема для окружности}
Центр $(0,0)$, $R=5$. Начальные значения: $x=0$, $y=5$, $d = 3 - 2\cdot5 = -7$.
Пошагово:
\begin{center}
\begin{tabular}{cccl}
\toprule
$x$ & $y$ & $d$ & Действие \\
\midrule
0 & 5 & $-7$  & $d<0 \Rightarrow d{+}{=}4\cdot0+6=6$ (стало $-1$) \\
1 & 5 & $-1$  & $d<0 \Rightarrow d{+}{=}10$ (стало $9$) \\
2 & 5 & $9$   & $d\ge0 \Rightarrow d{+}{=}4(2-5)+10=-2$ (стало $7$), $y{=}4$ \\
3 & 4 & $7$   & $d\ge0 \Rightarrow d{+}{=}4(3-4)+10=6$ (стало $13$), $y{=}3$ \\
\bottomrule
\end{tabular}
\end{center}

После применения симметрии по восьми октантам получаем $28$ уникальных точек:
\begin{multline*}
(0,5), (5,0), (0,-5), (-5,0),\ (1,5), (5,1), (5,-1), (1,-5),\\
(-1,-5), (-5,-1), (-5,1), (-1,5),\ (2,5), (5,2), (5,-2), (2,-5),\\
(-2,-5), (-5,-2), (-5,2), (-2,5),\ (3,4), (4,3), (4,-3), (3,-4),\\
(-3,-4), (-4,-3), (-4,3), (-3,4).
\end{multline*}

% ======================================
\section{Временные характеристики}

Замер времени осуществляется встроенной функцией \texttt{performance.now()}, которая обеспечивает точность до микросекунд. Для каждого алгоритма выполняется $N$ повторных построений (для отрезков $N = 20000$, для окружности $N = 5000$), после чего вычисляется среднее время одного построения. Результаты зависят от процессора и браузера; в таблице~\ref{tab:timing} приведены характерные значения для машины с процессором Intel Core i5 и браузером Chrome.

\begin{table}[H]
\centering
\caption{Усреднённое время построения одного примитива}
\label{tab:timing}
\begin{tabular}{lrrr}
\toprule
Алгоритм & Повторов & Среднее время, мс & Кол-во точек \\
\midrule
Пошаговый (отрезок)        & 20000 & $\sim 1{,}2\cdot10^{-3}$ & 10 \\
ЦДА                        & 20000 & $\sim 1{,}0\cdot10^{-3}$ & 10 \\
Брезенхем (отрезок)        & 20000 & $\sim 5{,}5\cdot10^{-4}$ & 10 \\
Брезенхем (окружность)     & 5000  & $\sim 4{,}0\cdot10^{-3}$ & 28 \\
\bottomrule
\end{tabular}
\end{table}

\noindent\textbf{Анализ.} Брезенхем для отрезка оказывается быстрее ЦДА и пошагового алгоритма примерно в 2 раза за счёт отсутствия операций деления и округления с плавающей точкой. Алгоритм для окружности выполняется дольше из-за большого числа операций \texttt{push} (до $8$ точек на итерацию), однако внутри цикла также использует только целочисленную арифметику. Полученные значения подтверждают теоретические оценки сложности: все четыре алгоритма имеют линейную сложность по длине примитива.

% ======================================
\section{Реализация дополнительных возможностей}

\subsection{Анимация}
При нажатии кнопки \textit{Анимация} выполняется пошаговое построение: точки выводятся с интервалом $\approx 80$~мс, а текущая строка таблицы растровых точек подсвечивается. Реализация основана на рекурсивном вызове \texttt{setTimeout}, что позволяет визуально наблюдать последовательность работы алгоритма.

\subsection{Масштабирование}
Регулятор масштаба (10–80 пикселей на единицу) позволяет исследовать поведение растеризации при разном разрешении растра. При изменении масштаба вся сцена перерисовывается без повторного вычисления массива точек.

\subsection{Замер времени}
Кнопка \textit{Замер времени} циклически вызывает каждый алгоритм и выводит таблицу с числом повторов, средним и суммарным временем. Это позволяет сравнить эффективность методов в реальном времени на конкретной машине.

% ======================================
\section{Заключение}

В ходе выполнения лабораторной работы были реализованы четыре базовых растровых алгоритма: пошаговый, ЦДА, Брезенхема для отрезка и Брезенхема для окружности. Все алгоритмы корректно строят растровые примитивы, что подтверждается совпадением результатов с ручными расчётами. Разработано веб-приложение с удобным интерфейсом, поддержкой сетки, осей координат, подписей, масштабирования, анимации и замера времени. 

Сравнение временных характеристик подтвердило теоретические ожидания: целочисленный алгоритм Брезенхема для отрезка работает примерно вдвое быстрее ЦДА и пошагового алгоритма. Пояснена схема привязки целочисленных координат к дискретной сетке через преобразование $s_x = W/2 + xs$, $s_y = H/2 - ys$.

Полученные навыки могут быть применены при разработке графических редакторов, систем компьютерной графики и визуализации данных.

% ======================================
\appendix
\section{Ключевые фрагменты кода}

\subsection{Алгоритм Брезенхема для отрезка}
\begin{lstlisting}[caption={bresenhamLine(x1, y1, x2, y2)}]
function bresenhamLine(x1, y1, x2, y2) {
  const pts = [];
  let x = x1, y = y1;
  const dx = Math.abs(x2 - x1);
  const dy = -Math.abs(y2 - y1);
  const sx = x1 < x2 ? 1 : -1;
  const sy = y1 < y2 ? 1 : -1;
  let err = dx + dy;
  while (true) {
    pts.push({ x, y });
    if (x === x2 && y === y2) break;
    const e2 = 2 * err;
    if (e2 >= dy) { err += dy; x += sx; }
    if (e2 <= dx) { err += dx; y += sy; }
  }
  return pts;
}
\end{lstlisting}

\subsection{Алгоритм Брезенхема для окружности}
\begin{lstlisting}[caption={bresenhamCircle(cx, cy, r)}]
function bresenhamCircle(cx, cy, r) {
  const pts = [];
  if (r < 0) return pts;
  let x = 0, y = r, d = 3 - 2 * r;
  function add8(px, py) {
    pts.push({x: cx+px, y: cy+py}, {x: cx+py, y: cy+px},
             {x: cx+py, y: cy-px}, {x: cx+px, y: cy-py},
             {x: cx-px, y: cy-py}, {x: cx-py, y: cy-px},
             {x: cx-py, y: cy+px}, {x: cx-px, y: cy+py});
  }
  while (x <= y) {
    add8(x, y);
    if (d < 0) d += 4 * x + 6;
    else { d += 4 * (x - y) + 10; y--; }
    x++;
  }
  return pts;
}
\end{lstlisting}

\subsection{Привязка координат к дискретной сетке}
\begin{lstlisting}[caption={worldToScreen(x, y)}]
function worldToScreen(x, y) {
  return {
    sx: W / 2 + x * state.scale,
    sy: H / 2 - y * state.scale
  };
}
\end{lstlisting}

\section*{Список использованных источников}
\begin{enumerate}[label={[\arabic*]}]
  \item Роджерс Д.~Ф. Алгоритмические основы машинной графики. — М.: Мир, 1989.
  \item Hearn D., Baker M.~P. Computer Graphics with OpenGL. — 4th ed. Pearson, 2011.
  \item Bresenham J.~E. Algorithm for computer control of a digital plotter // IBM Systems Journal, 1965.
\end{enumerate}

\end{document}